\section{Spectral estimates at a slowly enlarged edge scale}
\label{sec:spectral}

Our objective is to control the leading GOE eigenvalues and the first
two resolvent traces just above the edge on an event independent of
eigenvector orientation. These estimates enter both the Gaussian
conditioning argument in Section~\ref{sec:spherical} and the variational
deficit in Section~\ref{sec:orientation}. A fixed logarithmic enlargement
would lose the strongest quantifiers of the main theorems, so the
enlargement below may diverge as slowly as needed.

The proof has three parts. Tridiagonalization isolates independent noise;
a simultaneous quadratic-form estimate compares that noise with a
deterministic confinement energy; and a Hermite identity evaluates the
deterministic resolvent. Here $C,c>0$ are universal, while $C_M$ may also
depend on a fixed $M<\infty$.

Let $L_n\to\infty$ be any positive deterministic sequence with $L_n\le\log n$ for
all sufficiently large $n$, and set
\begin{equation}
 s_*=n^{-2/3},\qquad s_0=L_ns_*,\qquad
 h_n=L_n^{1/8},\qquad
 A_n=h_n+\sqrt{\log(2+L_n)},\qquad
 \eta_n=A_nL_n^{-3/4}.
 \label{eq:spectral-scales}
\end{equation}
In particular, $\eta_n\to0$.  Define
\[
 R(s)=\frac{2+s-\sqrt{s(4+s)}}2,\qquad s>0.
\]

\begin{proposition}[An eigenvalue-only good event]
\label{prop:spectral}
Let $W$ have the GOE normalization of Section~\ref{sec:model}, and write its
eigenvalues as $\lambda_1\ge\cdots\ge\lambda_n$.  There are events
$\mathcal G_n(L_n)$, measurable with respect to these eigenvalues alone,
such that
\[
 \Prob\bigl(\mathcal G_n(L_n)^c\bigr)
 \le C\exp(-cL_n^{1/4})
\]
for all sufficiently large $n$.  On the same events, for every fixed
$M<\infty$,
\begin{equation}
\begin{aligned}
 &\lambda_j=2+o(s_0)\quad(1\le j\le3),
 &&\max_i|\lambda_i|\le3,\\
 &\frac1n\sum_{i=1}^n\frac1{2+s-\lambda_i}
     =R(s)+o(\sqrt s),
 &&\frac1n\sum_{i=1}^n\frac1{(2+s-\lambda_i)^2}
     \le\frac{C_M}{\sqrt s}
 \quad(s_0\le s\le M).
\end{aligned}
\label{eq:spectral-estimates}
\end{equation}
All denominators in these formulas are positive.  The little-oh terms have
deterministic bounds, uniformly over the event: the first three eigenvalue
errors divided by $s_0$ are at most $C(L_n^{-1}+\eta_n)$, and the first
resolvent error divided by $\sqrt s$ is at most
$C_M(L_n^{-3/2}+\eta_n)$.  The event does not depend on $M$.
\end{proposition}

\subsection{Tridiagonal reduction and concentration of its entries}

The following reduction is the GOE tridiagonalization of
\cite[Theorem~2.1]{dumitriuEdelman2002}. We include the argument to fix both its
normalization and its independence properties.

\begin{lemma}[Tridiagonal representation]
\label{lem:tridiagonal}
The eigenvalues of $W$ have the same law as those of a symmetric
tridiagonal matrix $T$ whose entries
\[
 T_{ii}\sim N(0,2/n),\qquad
 T_{i,i+1}=T_{i+1,i}\sim\frac{\chi_{n-i}}{\sqrt n}
 \quad(1\le i<n)
\]
are independent.  Here $\chi_m$ is the Euclidean norm of $m$ independent
standard normal variables.
\end{lemma}

\begin{proof}
We expose one column at a time. The first diagonal entry, the rest of
the first column, and the trailing square block are independent; this
independence will persist after each rotation.  The column vector has $n-1$ independent $N(0,1/n)$ entries.
Choose an orthogonal transformation of the trailing coordinates that sends
this vector to its nonnegative length times the first coordinate vector.
Conjugation by this transformation produces the desired first
off-diagonal entry.

To justify the next exposure, observe that the trailing block has
density proportional to $\exp(-n\Tr B^2/4)$. Both this density and
Euclidean volume are invariant under orthogonal conjugation.  Conditional on the
column vector, the chosen transformation is fixed, so the conjugated
trailing block still has the same GOE density and its conditional law does
not depend on the column.  It is therefore independent of all entries
already exposed.  Repeat the construction inside that block, then inside
each successive trailing block.  At step $i$ the fresh column has $n-i$
independent $N(0,1/n)$ entries.  This proves both the stated distributions
and their joint independence.  The conjugations preserve eigenvalues.
\end{proof}

The reduction leaves Gaussian diagonal noise and centered chi noise
off the diagonal. Both satisfy a dimension-independent exponential
estimate. For the latter, the two estimates we need are
\begin{equation}
 \E\exp\left(t\frac{\chi_m-\E\chi_m}{\sqrt n}\right)
 \le \exp(Ct^2/n),\qquad
 0\le\sqrt m-\E\chi_m\le\frac C{\sqrt m}.
 \label{eq:chi-entry-bounds}
\end{equation}
Here and below a moment-generating bound is asserted for every real $t$.
For completeness, the Gaussian-square moment-generating function and
Chernoff's inequality give
\[
 \Prob\{\chi_m^2\ge ma\}\le(ae^{1-a})^{m/2}
 \quad(a>1),
\]
with the same right-hand side for the lower-tail event when $0<a<1$.
Since
\[
 a-1-\log a\ge(\sqrt a-1)^2,
\]
it follows that
$\Prob\{|\chi_m-\sqrt m|\ge r\}\le2e^{-r^2/2}$ for $r\ge0$.
This tail bound gives $\Var(\chi_m)\le C$ and, by Jensen's inequality,
\[
 \sqrt m-\E\chi_m
 =\frac{\Var(\chi_m)}{\sqrt m+\E\chi_m}\le\frac C{\sqrt m}.
\]
To prove the centered exponential bound without an extra prefactor, let
$X,X'$ be independent copies of $\chi_m$ and put $Z=X-X'$.  Then $Z$ is
symmetric and
\[
 \Prob\{|Z|\ge r\}\le4e^{-r^2/8},\qquad
 \E Z^{2k}\le4\cdot8^k k!\quad(k\ge1).
\]
The moment bound follows by integrating the tail.  Jensen's inequality
conditional on $X$, followed by the power series of $\cosh$, gives
\[
\begin{aligned}
 \E e^{t(X-\E X)}
 &\le\E e^{tZ}=\E\cosh(tZ)\\
 &\le1+4\sum_{k\ge1}\frac{8^k k!t^{2k}}{(2k)!}
 \le1+4(e^{8t^2}-1)\le e^{32t^2}.
\end{aligned}
\]
We used $(2k)!\ge(k!)^2$ and $(1+u)^4\ge1+4u$ for $u\ge0$.
The Gaussian tail ensures exponential integrability, and Tonelli's theorem
also justifies the nonnegative $\cosh$ series.  Replacing $t$ by
$t/\sqrt n$ proves~\eqref{eq:chi-entry-bounds}.  The centered diagonal
entries satisfy the same exponential estimate directly.

\subsection{A simultaneous quadratic-form estimate}

The quadratic-form estimate follows the approach used in soft-edge analysis for
tridiagonal ensembles. Ram\'irez, Rider, and Vir\'ag
\cite{ramirezRiderVirag2011} developed a variational approach involving
discrete gradient and confinement energies. The block summation-by-parts
estimate of Ledoux and Rider \cite[Proposition~7 and Lemma~8]{ledouxRider2010}
is particularly close to the noise control used here. Their proof of
Theorem~1 (Right-Tail), on page~1329 immediately before Lemma~9, gives
the corresponding substitution for off-diagonal weights. We arrange that control simultaneously
over energy scales and spatial shells, then use a deterministic Hermite
resolvent identity. This provides the first and second resolvent bounds
needed here for every slowly diverging edge enlargement.

Let $D$ be the deterministic tridiagonal matrix with zero diagonal and
off-diagonal entries
\[
 c_i=\sqrt{1-i/n}\quad(1\le i<n),\qquad E=2I-D.
\]
For later boundary formulas set $c_0=c_n=0$.

\begin{lemma}[Noise controlled by deterministic energy]
\label{lem:form-bound}
There is a universal constant $C_0$ such that, with
\[
 \Phi_n(x)=C_0n^{-1/2}A_n(x+s_0)^{1/4}
             \sqrt{1+\log(1+x/s_0)},\qquad x\ge0,
\]
the matrix $H=T-D$ satisfies, simultaneously for every unit vector
$u\in\R^n$,
\begin{equation}
 |u^THu|\le\Phi_n(u^TEu)
 \label{eq:form-bound}
\end{equation}
on an event $\mathcal F_n$ with
$\Prob(\mathcal F_n^c)\le C\exp(-cL_n^{1/4})$ for all sufficiently large $n$.
Moreover,
\begin{equation}
 \sup_{x\ge0}|\Phi_n'(x)|\le C\eta_n,
 \qquad \Phi_n(0)\le C\eta_ns_0.
 \label{eq:phi-smallness}
\end{equation}
\end{lemma}

\begin{proof}
We prove the lemma in three steps: the deterministic energy controls
where a vector can concentrate, one noise event controls block prefixes
at every relevant scale, and summation by parts combines these controls.

\smallskip\noindent\emph{Energy controls spatial mass and gradients.}
For a unit vector $u$, write $e=u^TEu$. Expanding gives
\begin{equation}
 e=\sum_{i<n}c_i(u_i-u_{i+1})^2
   +\sum_{i=1}^n(2-c_i-c_{i-1})u_i^2.
 \label{eq:deterministic-energy}
\end{equation}
The row sums of $D$ are at most two, so $0\le E\le4I$.  The coefficients
$2-c_i-c_{i-1}$ are at least $i/(2n)$: for interior indices this follows
from $1-\sqrt{1-t}\ge t/2$, and the boundary indices satisfy the same
inequality directly.  Also $c_i\ge1/\sqrt2$ when $i\le n/2$.
Thus~\eqref{eq:deterministic-energy} implies
\begin{equation}
 \sum_i\frac{i}{n}u_i^2\le2e,\qquad
 \sum_{i>n/2}u_i^2+
 \sum_{i<n}(u_i-u_{i+1})^2\le Ce.
 \label{eq:energy-tail-gradient}
\end{equation}
For the second inequality, control gradients in the first half by the
weighted gradient term, and use
$(u_i-u_{i+1})^2\le2u_i^2+2u_{i+1}^2$ for the remaining edges.

\smallskip\noindent\emph{One event for every energy.}
The vector will be selected after the matrix has been realized, so the
probability estimate must be simultaneous. We first prepare blocks at
every dyadic energy scale. Set
\[
 \sigma_j=2^js_*,\qquad
 0\le j\le J_n:=\left\lceil\log_2(4/s_*)\right\rceil,
 \qquad d_j=\max\{1,\lfloor\sigma_j^{-1/2}\rfloor\}.
\]
For each of the two independent-entry sequences (diagonal entries and
centered off-diagonal entries), partition its index set into consecutive
blocks of lengths between $d_j$ and $2d_j$.  Such a partition is obtained
by merging a last short block with its predecessor.  Assign label $r=0$
when the first index of a block is at most $2n\sigma_j$; otherwise assign
$r\ge1$ when that index is in
$(2^rn\sigma_j,2^{r+1}n\sigma_j]$.

If $\xi_i$ is either centered entry sequence and a block has length
$m\le2d_j$, its partial sums $S_k$ obey
\[
 \Prob\left\{\max_{k\le m}|S_k|\ge a\right\}
 \le2\exp\left(-\frac{na^2}{C d_j}\right).
\]
Indeed, $\exp(tS_k-Ckt^2/n)$ is a nonnegative supermartingale by
\eqref{eq:chi-entry-bounds}.  Stopping it at its first crossing of $a$
bounds the one-sided probability by
$\exp(-ta+Cmt^2/n)$; optimizing $t>0$ and repeating for $-S_k$ proves
the displayed estimate.

At a fixed pair $(j,r)$, the number of blocks is at most
\[
 C\left(1+\frac{2^rn\sigma_j}{d_j}\right)
 \le C2^rn\sigma_j^{3/2}=C2^{r+3j/2}.
\]
For $\sigma_j\le1$, use $d_j\ge\sigma_j^{-1/2}/2$; for
$\sigma_j>1$, use $d_j=1$.  The additive constant is absorbed because
$n\sigma_j^{3/2}\ge1$.  Consequently, for a sufficiently large fixed
constant $K$, there is an event $\mathcal F_n$ on which every block of
both sequences has all its centered partial sums bounded by
\begin{equation}
 K\sqrt{d_j/n}\,(h_n+\sqrt{j+r}),
 \label{eq:block-prefix}
\end{equation}
and
\[
\begin{aligned}
 \Prob(\mathcal F_n^c)
 &\le C\sum_{j,r\ge0}2^{r+3j/2}
       e^{-cK^2(h_n+\sqrt{j+r})^2}\\
 &\le C e^{-c'h_n^2}=C e^{-c'L_n^{1/4}}.
\end{aligned}
\]
Here $K$ is chosen so that the factors in $j$ and $r$ form summable
geometric series.  This bound does not acquire a factor from the number
of scales.

\smallskip\noindent\emph{From block prefixes to quadratic forms.}
We may now fix the noise event and choose an arbitrary unit vector.
Select the first $j$ for which $e\le\sigma_j$. Since $e\le4$, this
index exists, and
\[
 \sigma_j\le2(e+s_*).
\]
For the weights $w_i=u_i^2$ or $w_i=u_iu_{i+1}$, the sum of their
absolute values over all blocks with label $r$ is at most $C2^{-r}$.
For $r\ge1$, every coordinate used is beyond $2^rn\sigma_j$, and
\eqref{eq:energy-tail-gradient} bounds its total squared mass by
$2e/(2^r\sigma_j)\le2^{1-r}$.  The off-diagonal weights obey the same
bound by $|u_iu_{i+1}|\le(u_i^2+u_{i+1}^2)/2$.  For $r=0$, use the
unit norm.

The summed variation of these weights within those blocks is at most
$C2^{-r/2}\sqrt e$.  To verify this assertion, use
\[
 u_{i+1}^2-u_i^2=(u_{i+1}-u_i)(u_{i+1}+u_i)
\]
for diagonal weights and
\[
 u_{i+1}u_{i+2}-u_iu_{i+1}
 =u_{i+1}\bigl[(u_{i+2}-u_{i+1})+(u_{i+1}-u_i)\bigr]
\]
for off-diagonal weights.  Cauchy--Schwarz combines the preceding tail
mass bound with the total gradient bound in
\eqref{eq:energy-tail-gradient}; each coordinate or edge is counted at
most a fixed number of times.

For clarity, the summation-by-parts step can be made blockwise.  Suppose
all centered entry prefixes in a block are bounded by $B$, and let
$\bar\xi$ be the entry average on that block.  Then
$|\bar\xi|\le B/d_j$.  The prefixes of $\xi_i-\bar\xi$ are at most
$2B$ and vanish at the block end.  Therefore
\[
 \left|\sum_{i\text{ in block}}\xi_iw_i\right|
 \le\frac B{d_j}\sum_{i\text{ in block}}|w_i|
      +2B\sum_{\substack{i,i+1\text{ in block}}}|w_{i+1}-w_i|.
\]
Using~\eqref{eq:block-prefix} and summing by labels proves that either
centered weighted sum is bounded by
\begin{align*}
 &C n^{-1/2}\sum_{r\ge0}(h_n+\sqrt{j+r})
 \left(d_j^{-1/2}2^{-r}+d_j^{1/2}2^{-r/2}\sqrt e\right)\\
 &\hspace{35mm}\le C(h_n+\sqrt j)n^{-1/2}\sigma_j^{1/4}.
\end{align*}
For $\sigma_j\le1$, the last step uses
$d_j\asymp\sigma_j^{-1/2}$ and $e\le\sigma_j$.  For
$1<\sigma_j<8$, it follows from $d_j=1$ with a changed universal
constant.  The series in $r$ converges geometrically, including its
$\sqrt r$ factor.

The centered entries have now been controlled uniformly in $u$.
The only remaining part of $T-D$ is the deterministic error in the
off-diagonal means. By
\eqref{eq:chi-entry-bounds}, their difference from $c_i$ is at most
$C/(\sqrt n\sqrt{n-i})$.  It is at most $C/n$ for $i\le n/2$ and at
most $C/\sqrt n$ thereafter.  Thus their quadratic-form contribution is
at most
\[
 C/n+Ce/\sqrt n,
\]
using again the tail mass bound.  Finally,
\[
 j\le C\log(2+e/s_*)
 \le C\bigl[\log(2+L_n)+\log(1+e/s_0)\bigr].
\]
Together with $\sigma_j\le2(e+s_*)$ and $s_*\le s_0$, this converts
the centered estimate into~\eqref{eq:form-bound}.  The mean errors are
absorbed into the same bound because $0\le e\le4$ and $L_n\ge1$ for
large $n$.  Both sequences, and the factor two from off-diagonal entries,
are absorbed by choosing $C_0$ once and for all.

Finally, the scalar comparison must be stable under changing the
energy. Put $q(x)=1+\log(1+x/s_0)$. Direct differentiation gives
\[
 \Phi_n'(x)=C_0n^{-1/2}A_n(x+s_0)^{-3/4}
 \left(\frac{\sqrt{q(x)}}4+\frac1{2\sqrt{q(x)}}\right).
\]
Since $y^{-3/4}\sqrt{1+\log y}$ is bounded for $y\ge1$, its supremum
is at most $Cn^{-1/2}A_ns_0^{-3/4}=C\eta_n$.
The equality $\Phi_n(0)/s_0=C_0\eta_n$ proves the other assertion in
\eqref{eq:phi-smallness}.
\end{proof}

\subsection{The deterministic resolvent}
The noise bound reduces the random trace problem to $D$. Its entries
were chosen so that the characteristic polynomials satisfy a Hermite
recurrence. The resulting scalar differential equation will control both
the resolvent and the number of eigenvalues near the edge.

\begin{lemma}[Hermite resolvent estimates]
\label{lem:deterministic-resolvent}
Let $e_1\le\cdots\le e_n$ be the eigenvalues of $E=2I-D$, and set
\[
 R_D(s)=\frac1n\Tr(s+E)^{-1},\qquad
 T_D(s)=\frac1n\Tr(s+E)^{-2}=-R_D'(s).
\]
For every fixed $M<\infty$ and all sufficiently large $n$, uniformly for
$s_0\le s\le M$,
\begin{equation}
 T_D(s)\le C_Ms^{-1/2},\qquad
 0\le R(s)-R_D(s)\le\frac{C_M}{ns}.
 \label{eq:deterministic-traces}
\end{equation}
Moreover,
\begin{equation}
 N(t):=\#\{i:e_i\le t\}\le Cnt^{3/2}
 \qquad(t\ge s_0).
 \label{eq:deterministic-counting}
\end{equation}
\end{lemma}

\begin{proof}
We first derive the differential equation and locate its relevant
quadratic root, then absorb its finite-dimensional error at $s\ge s_0$.
Reverse the coordinates of $D$. Its off-diagonal entries become
$\sqrt{k/n}$, $1\le k<n$, so the characteristic polynomials of its
successive principal minors satisfy
\[
 P_0(z)=1,\quad P_1(z)=z,\quad
 P_{k+1}(z)=zP_k(z)-(k/n)P_{k-1}(z).
\]
Induction gives $P_k'=kP_{k-1}$: differentiation of the recurrence and
the recurrence one step earlier show that
$P_{k+1}'=P_k+kP_k$.  Consequently
\[
 P_n''-nzP_n'+n^2P_n=0.
\]
Since $P_n$ is the characteristic polynomial of $D$,
$R_D(s)=P_n'(2+s)/(nP_n(2+s))$.  The differential equation yields
\begin{equation}
 R_D(s)^2-(2+s)R_D(s)+1=\frac{T_D(s)}n>0.
 \label{eq:hermite-riccati}
\end{equation}
For every $s>0$, $R_D(s)$ is finite because $E\ge0$.
The two roots of the quadratic on the left are $R(s)$ and
\[
 R_+(s)=\frac{2+s+\sqrt{s(4+s)}}2.
\]
At large $s$, $R_D(s)\to0$ whereas $R_+(s)\to\infty$.  Positivity
in~\eqref{eq:hermite-riccati} thus places $R_D$ below the smaller root
there, and continuity prevents it from crossing either root anywhere on
$(0,\infty)$.  If $\delta(s)=R(s)-R_D(s)$, factoring the quadratic gives
\begin{equation}
 0\le\delta(s)\le
 \frac{T_D(s)}{n\sqrt{s(4+s)}}.
 \label{eq:root-difference}
\end{equation}

Termwise comparison of resolvents shows that
\[
 T_D(s)\le\frac2s\bigl[R_D(s/2)-R_D(s)\bigr]
 \le\frac2s\bigl[R(s/2)-R(s)+\delta(s)\bigr].
\]
Differentiating the explicit expression for $R$ gives
\[
 0<-R'(s)=\frac{2}{\sqrt{s(4+s)}\,
          (2+s+\sqrt{s(4+s)})}\le\frac1{2\sqrt s}.
\]
Integrating from $s/2$ to $s$ therefore gives
$R(s/2)-R(s)\le\sqrt s-\sqrt{s/2}\le\sqrt s$.  Hence
\[
 T_D(s)\le\frac C{\sqrt s}
       +\frac{2T_D(s)}{ns\sqrt{s(4+s)}}.
\]
For $s\ge s_0$, the coefficient of the last term is at most
$(ns_0^{3/2})^{-1}=L_n^{-3/2}\to0$.  Absorbing it proves the bound on
$T_D$, and~\eqref{eq:root-difference} then proves the bound on
$\delta$.  Notice that $R_D(s/2)\le R(s/2)$ was established for every
positive $s/2$; it does not require $s/2\ge s_0$.

The second-trace estimate also gives the counting estimate needed for
the later dyadic transfer. For $s_0\le t\le1$, each $e_i\le t$
contributes at least
$(2t)^{-2}$ to $nT_D(t)$.  Thus
$N(t)\le4nt^2T_D(t)\le Cnt^{3/2}$.  For $t\ge1$, the same bound
follows from $N(t)\le n$.
\end{proof}

\subsection{Eigenvalue comparison and transfer of the traces}

\begin{proof}[Proof of Proposition~\ref{prop:spectral}]
The previous lemmas supply the noise control and the reference traces.
We now transfer them to an event defined entirely by the eigenvalues.
There are three tasks: compare ordered eigenvalues, locate the top three,
and compare the two traces.

\smallskip\noindent\emph{Ordered eigenvalues.}
Work first in the tridiagonal realization. On $\mathcal F_n$,
\eqref{eq:form-bound} bounds the Rayleigh quotient of $E-H=2I-T$ between
\[
 f_-(e):=e-\Phi_n(e)\quad\hbox{and}\quad f_+(e):=e+\Phi_n(e).
\]
By~\eqref{eq:phi-smallness}, both functions are increasing for large $n$.
The span of the first $i$ eigenvectors of $E$ has every unit-vector
energy at most $e_i$, so the min--max principle gives
$2-\lambda_i\le f_+(e_i)$.  On the span of eigenvectors $i$ through
$n$, every energy is at least $e_i$; the complementary max--min
principle gives $2-\lambda_i\ge f_-(e_i)$.  Therefore
\begin{equation}
 |2-\lambda_i-e_i|\le\Phi_n(e_i),\qquad 1\le i\le n.
 \label{eq:eigenvalue-comparison}
\end{equation}
This argument uses monotonicity of scalar functions of Rayleigh
quotients, and does not require an operator-monotonicity assertion.

Define $\mathcal G_n(L_n)$ to be precisely the event
\eqref{eq:eigenvalue-comparison}, with the deterministic $e_i$ and
$\Phi_n$ constructed above.  It depends only on the ordered eigenvalues.
The tridiagonal event $\mathcal F_n$ implies it; hence
Lemma~\ref{lem:tridiagonal} transfers its probability bound to $W$.
All remaining conclusions are deterministic on $\mathcal G_n(L_n)$.

\smallskip\noindent\emph{The first three directions.}
The conditioning argument will need three large coordinate variances.
To locate the corresponding eigenvalues, take
$m=\lfloor n^{1/3}\rfloor$ and the span of the first three discrete
Dirichlet sine modes on $\{1,\ldots,m\}$, extended by zero.
The Dirichlet Laplacian on this segment has eigenvalues
$4\sin^2(k\pi/[2(m+1)])$, so its quadratic form on this span is at most
$C/m^2$.  On vectors supported on the segment, the form difference
between $E$ and that Laplacian has absolute value at most
\[
 2\max_{1\le i<m}|1-c_i|\sum_{i<m}|u_iu_{i+1}|
 \le Cm/n.
\]
The min--max principle therefore gives
$0\le e_1,e_2,e_3\le C(m^{-2}+m/n)\le Cn^{-2/3}$.
Equations~\eqref{eq:phi-smallness} and
\eqref{eq:eigenvalue-comparison} imply
\[
 \max_{j\le3}\frac{|\lambda_j-2|}{s_0}
 \le C(L_n^{-1}+\eta_n).
\]
Also every $e_i\in[0,4]$, so
$\max_i|\lambda_i|\le2+\Phi_n(4)$.  The assumption $L_n\le\log n$
gives $\Phi_n(4)=o(1)$, proving the stated bound by three.

\smallskip\noindent\emph{Resolvent traces.}
The eigenvalue comparison also controls every denominator: for every
$s\ge s_0$ and every $e\in[0,4]$,
\begin{equation}
 \frac{\Phi_n(e)}{s+e}
 \le\frac{\Phi_n(0)+C\eta_ne}{s+e}
 \le C\eta_n.
 \label{eq:denominator-comparison}
\end{equation}
Thus~\eqref{eq:eigenvalue-comparison} makes
$2+s-\lambda_i$ positive and equal to $(s+e_i)(1+O(\eta_n))$,
uniformly in both $i$ and $s\ge s_0$.
The normalized second trace is consequently at most a constant times
$T_D(s)$, proving its bound in~\eqref{eq:spectral-estimates}.

For the first trace, subtract term by term and use
$(e_i+s_0)^{1/4}\le(s+e_i)^{1/4}$ to obtain
\begin{align}
 &\left|\frac1n\sum_i\frac1{2+s-\lambda_i}-R_D(s)\right|\notag\\
 &\qquad\le Cn^{-1/2}A_n\frac1n\sum_i
 (s+e_i)^{-7/4}\sqrt{1+\log(1+e_i/s_0)}.
 \label{eq:trace-perturbation}
\end{align}
We spell out the dyadic sum because its uniformity near $s_0$ is useful.
The range $e_i\le s$ contributes at most
\[
 Cs^{-7/4}\frac{N(s)}n\sqrt{1+\log(1+s/s_0)}
 \le Cs^{-1/4}\sqrt{1+\log(1+s/s_0)}.
\]
For $k\ge0$, the shell $2^ks<e_i\le2^{k+1}s$ contributes at most
\[
 Cs^{-1/4}2^{-k/4}
 \sqrt{1+\log(1+s/s_0)+(k+1)\log2},
\]
by~\eqref{eq:deterministic-counting}.  This counting estimate remains
valid above one, so all shells are covered.  The convergent geometric
series, with its additional square-root factor, bounds the sum in
\eqref{eq:trace-perturbation} by
\[
 C_Ms^{-1/4}\sqrt{1+\log(1+s/s_0)}.
\]
Dividing the resulting trace error by $\sqrt s$ and writing $r=s/s_0$
gives
\[
 C_M\eta_n r^{-3/4}\sqrt{1+\log(1+r)}\le C_M\eta_n
 \qquad(r\ge1).
\]
Finally, Lemma~\ref{lem:deterministic-resolvent} gives
\[
 \frac{|R_D(s)-R(s)|}{\sqrt s}
 \le\frac{C_M}{ns^{3/2}}\le C_ML_n^{-3/2}.
\]
This proves the quantitative first-trace error and completes all of
\eqref{eq:spectral-estimates}.  The event was defined before $M$ was
chosen, and each displayed error is deterministic on that event.
Consequently every deterministic sequence of spectra in
$\mathcal G_n(L_n)$ satisfies the estimates for every fixed finite $M$,
without any condition on eigenvector orientation.
\end{proof}
